\documentclass[reqno, 12pt]{amsart}
\usepackage[brazil]{babel}
\usepackage[utf8]{inputenc}
\usepackage[T1]{fontenc}
\pagestyle{plain}
\usepackage{amsfonts,amssymb,amsmath,textcomp,amsthm}
\setlength{\parindent}{1cm}
\usepackage{ae}
\date{}

\usepackage[]{hyperref}

\newenvironment{packed_itemize}{
\begin{itemize}
  \setlength{\itemsep}{2pt}
  \setlength{\parskip}{0pt}
  \setlength{\parsep}{0pt}
}{\end{itemize}}

\usepackage{geometry}
\geometry{left=23.5mm,right=23.5mm, top=13mm, bottom=25mm}

\usepackage{setspace}
\onehalfspacing
\renewcommand{\baselinestretch}{1.2}

\begin{document}

\begin{center}
\vspace{.2cm}
\textbf{\Large Plano de Trabalho para mudança de regimento}\\
\vspace{.4cm}
\end{center}

\noindent \textit{Aluno}: {\bf Afonso Sant'Anna}\\
\noindent \textit{Orientador}: \textbf{Prof. Dr. Yoshiharu Kohayakawa}

\section{Descrição do Plano}

Atendendo à deliberação da CCP de 13/02/2026, que condiciona a migração ao novo regulamento à apresentação de plano de trabalho com cronograma semestral de disciplinas, encaminho o presente documento em conformidade com o item IV.4 da Resolução CoPGr n\textsuperscript{o} 8877, de 11 de novembro de 2025.

Declaro que meu orientador está ciente e de acordo com este plano.

O item IV.4 estabelece que os estudantes devem cumprir créditos em disciplinas formativas, assegurando formação sólida na área, incluindo ao menos uma disciplina do núcleo de Computação Teórica e uma do núcleo de Computação Aplicada, além de disciplinas adicionais relacionadas à tese ou pertinentes à formação especializada.

\subsection{ Núcleo de disciplinas formativas em Teoria da Computação e disciplinas adicionais relacionadas à tese}

\textbf{Durante o Doutorado:}
\begin{packed_itemize}
  \item MAC5711 -- Análise de Algoritmos (obrigatória)
\end{packed_itemize}

\textbf{Durante o Mestrado:}
\begin{packed_itemize}
  \item MAC5770 -- Introdução à Teoria dos Grafos (obrigatória)
  \item MAC4722 -- Linguagens, Autômatos e Computabilidade (obrigatória)
\end{packed_itemize}

\textbf{Disciplinas adicionais relacionadas à área de pesquisa (Combinatória) ou pertinentes à formação especializada:}

\textbf{Durante o Doutorado:}
\begin{packed_itemize}
  \item MAT6707 -- Métodos Algébricos em Combinatória
  \item MAC5005 -- Fundamentos Matemáticos do Aprendizado Sequencial e Online (em andamento)
  \item MAC6962 -- Tópicos Matemáticos para Computação Contemporânea (a ser cursada no 2º semestre de 2026)
\end{packed_itemize}

\textbf{Durante o Mestrado:}
\begin{packed_itemize}
  \item MAC5922 -- Combinatória I
  \item MAC5923 -- Combinatória II
  \item MAC6918 -- Tópicos na Teoria Algébrica dos Grafos
  \item MAC6989 -- Seminários em Combinatória Extremal, Probabilística e Aditiva III
  \item MAC6990 -- Seminários em Combinatória Extremal, Probabilística e Aditiva IV
  \item MAC6953 -- Seminários em Combinatória Extremal, Probabilística e Aditiva I
  \item MAC6954 -- Seminários em Combinatória Extremal, Probabilística e Aditiva II
\end{packed_itemize}

\subsection{ Núcleo de disciplinas formativas em Computação Aplicada} 

\textbf{Durante o Mestrado:}
\begin{packed_itemize}
  \item MAC5786 -- Princípios de Interação Humano--Computador (obrigatória)
\end{packed_itemize}

Dessa forma, no âmbito do Doutorado, as disciplinas MAC5711 (formativa — núcleo de Computação Teórica), MAT6707 (disciplina da área de pesquisa) e MAC6962 (disciplina da área de pesquisa, a ser cursada no segundo semestre de 2026) totalizarão ao menos 26 créditos, atendendo ao mínimo exigido pelo regulamento para o curso de Doutorado.

Assim, verifico que me encontro em conformidade com o item IV.4 do regulamento, possuindo formação no núcleo de Computação Teórica, no núcleo de Computação Aplicada e disciplinas adicionais diretamente relacionadas à minha área de pesquisa.

\subsection{Proficiência em Língua Inglesa}

Minha matrícula no Doutorado ocorreu em julho de 2025. Encontro-me dentro do prazo regulamentar para comprovação de proficiência em língua inglesa e realizarei o exame no semestre corrente.

\subsection{Exame de Qualificação}

Nos termos do regulamento, a inscrição no Exame de Qualificação deve ocorrer no prazo de até 24 meses após a matrícula inicial no Doutorado. Considerando minha data de ingresso, ainda disponho de prazo regulamentar suficiente para a realização do exame, prevista para o segundo semestre de 2026, em consonância com o cronograma de pesquisa.
\subsection{Cronograma de Pesquisa}

No que tange ao cronograma de pesquisa para o período remanescente do Doutorado, pretendo intensificar as investigações na área de Teoria de Ramsey, com foco inicial em problemas probabilísticos relacionados a resultados anti-Ramsey, bem como em problemas de Ramsey canônico.
 
Entre os dias 2 e 14 de fevereiro de 2026, participei do evento ``1° Friends in Brazilian Extremal Combinatorics'', realizado no IMPA (Rio de Janeiro, Brasil), durante o qual integrei sessões de trabalho colaborativo em problema relacionado a colorações canônicas em Teoria de Ramsey, em conjunto com Walner Mendonça, Guilherme Mota e Hugo Vicente. As discussões realizadas nesse período originaram desdobramentos que pretendo aprofundar ao longo do Doutorado.

Com vistas ao aprofundamento dessa linha de investigação, será encaminhado ainda neste semestre pedido formal de coorientação do Prof. Dr. Guilherme Mota, visando ao desenvolvimento conjunto de problemas na área de Teoria de Ramsey. Adicionalmente, planejo elaborar e submeter projeto de bolsa à FAPESP para o período remanescente do Doutorado.

Solicito, portanto, a migração para o novo regulamento do Programa de Pós-Graduação em Ciência da Computação, estabelecido pela Resolução CoPGr n\textsuperscript{o} 8877, de 11 de novembro de 2025.

\end{document}

%%% Local Variables:
%%% mode: latex
%%% eval: (auto-fill-mode t)
%%% eval: (LaTeX-math-mode t)
%%% eval: (flyspell-mode t)
%%% TeX-master: t
%%% End:
